% ===============================================================
% appendix.tex  --  Appendices (excluded from 10-page main limit)
% ===============================================================
\clearpage
\onecolumn
\appendix
\renewcommand{\thesection}{\Alph{section}}
\titleformat{\section}
  {\normalsize\bfseries\centering}
  {Appendix \thesection.}{0.5em}{}

% ---------------------------------------------------------------
\section{Full Evaluation Results}
\label{app:eval-table}

\renewcommand{\arraystretch}{1.15}
\begin{longtable}{@{}p{0.7cm} p{4.3cm} p{2.5cm} c c c c p{4.2cm}@{}}
\caption{Full evaluation results across all 15 questions and 3 systems.
Generation in the batch harness used the stand-in generator described in
Section~\ref{sec:evaluation}; retrieval scores are model-independent.
Cor.=Correctness, Grd.=Grounding, Ret.=Retrieval relevance,
Use.=Usefulness. ``--'' = not applicable or generation failed.}
\label{tab:full-eval}\\
\toprule
ID & Question & System & Cor. & Grd. & Ret. & Use. & Notes \\
\midrule
\endfirsthead
\multicolumn{8}{c}{\footnotesize\itshape Table~\ref{tab:full-eval} continued from previous page}\\
\toprule
ID & Question & System & Cor. & Grd. & Ret. & Use. & Notes \\
\midrule
\endhead
\midrule \multicolumn{8}{r}{\footnotesize\itshape continued on next page}\\
\endfoot
\bottomrule
\endlastfoot
Q01 & Who is Hamlet? & Baseline   & 3 & -- & 0 & 3 & Prompt repeated in output \\
    &                & RAG (scene)& 4 & 3  & 5 & 4 & Best result for this question \\
    &                & Enh.\ RAG  & 3 & 2  & 5 & 3 & Retrieved correct utterances \\
\midrule
Q02 & Role of Lady Macbeth? & Baseline   & 1 & -- & 0 & 1 & Completely hallucinated \\
    &                       & RAG (scene)& 3 & 2  & 5 & 3 & Retrieved murder scene \\
    &                       & Enh.\ RAG  & 3 & 2  & 5 & 3 & Similar to scene-level \\
\midrule
Q03 & Montagues vs.\ Capulets? & Baseline   & 1 & -- & 0 & 1 & No context available \\
    &                          & RAG (scene)& -- & -- & -- & -- & Generation failed \\
    &                          & Enh.\ RAG  & 3 & 2  & 5 & 3 & Utterance retrieval worked \\
\midrule
Q04 & Why does Macbeth kill Duncan? & Baseline   & 3 & -- & 0 & 3 & Generic answer \\
    &                               & RAG (scene)& 3 & 2  & 5 & 3 & Retrieved relevant scenes \\
    &                               & Enh.\ RAG  & 3 & 2  & 5 & 3 & Similar performance \\
\midrule
Q05 & Why does Hamlet delay revenge? & Baseline   & 3 & -- & 0 & 3 & Generic answer \\
    &                                & RAG (scene)& 3 & 2  & 5 & 3 & Retrieval correct \\
    &                                & Enh.\ RAG  & 2 & 1  & 5 & 2 & Weaker grounding \\
\midrule
Q06 & What happens to Ophelia? & Baseline   & 3 & -- & 0 & 3 & Generic \\
    &                          & RAG (scene)& 3 & 2  & 5 & 3 & Retrieved Ophelia scenes \\
    &                          & Enh.\ RAG  & 3 & 2  & 5 & 3 & Similar \\
\midrule
Q07 & R\&J relationship development? & Baseline   & 1 & -- & 0 & 1 & Hallucinated \\
    &                                & RAG (scene)& 3 & 2  & 5 & 3 & Retrieved key scenes \\
    &                                & Enh.\ RAG  & 3 & 2  & 5 & 3 & Similar \\
\midrule
Q08 & ``Wherefore art thou Romeo''? & Baseline   & 1 & -- & 0 & 1 & Missed the meaning \\
    &                               & RAG (scene)& 3 & 2  & 5 & 3 & Retrieved balcony scene \\
    &                               & Enh.\ RAG  & 3 & 2  & 5 & 3 & Similar \\
\midrule
Q09 & Ambition: Macbeth vs.\ Lady M? & Baseline   & 3 & -- & 0 & 3 & Generic \\
    &                                & RAG (scene)& 3 & 2  & 5 & 3 & Both characters covered \\
    &                                & Enh.\ RAG  & 3 & 2  & 5 & 3 & Similar \\
\midrule
Q10 & Role of the ghost? & Baseline   & 3 & -- & 0 & 3 & Generic \\
    &                    & RAG (scene)& 3 & 2  & 5 & 3 & Retrieved ghost scenes \\
    &                    & Enh.\ RAG  & 1 & 1  & 5 & 1 & Generation failed badly \\
\midrule
Q11 & Why does Juliet fake death? & Baseline   & 3 & -- & 0 & 3 & Generic \\
    &                             & RAG (scene)& 3 & 2  & 5 & 3 & Retrieved relevant scene \\
    &                             & Enh.\ RAG  & 3 & 2  & 5 & 3 & Similar \\
\midrule
Q12 & How does Macbeth change? & Baseline   & 3 & -- & 0 & 3 & Generic \\
    &                          & RAG (scene)& -- & -- & -- & -- & Generation failed \\
    &                          & Enh.\ RAG  & 3 & 2  & 5 & 3 & Multi-scene handled \\
\midrule
Q13 & Juliet stylised response & Baseline   & 3 & -- & 0 & 3 & No Shakespearean style \\
    &                          & RAG (scene)& 3 & 2  & 5 & 3 & Retrieved balcony scene \\
    &                          & Enh.\ RAG  & 3 & 2  & 5 & 3 & Style still weak \\
\midrule
Q14 & Significance of three witches? & Baseline   & 3 & -- & 0 & 3 & Generic \\
    &                                & RAG (scene)& 1 & 1  & 5 & 1 & Generation failed \\
    &                                & Enh.\ RAG  & 3 & 2  & 5 & 3 & Better than scene-level \\
\midrule
Q15 & Who is Mercutio? & Baseline   & 3 & -- & 0 & 3 & Generic \\
    &                  & RAG (scene)& 3 & 2  & 5 & 3 & Retrieved relevant scene \\
    &                  & Enh.\ RAG  & 3 & 2  & 5 & 3 & Similar \\
\end{longtable}

% ---------------------------------------------------------------
\section{Representative Prompts}
\label{app:prompts}

\noindent\textbf{RAG system prompt} (\texttt{prompts/system\_prompt.txt}),
used for \texttt{qa} and \texttt{concept} modes:
\begin{lstlisting}
You are a Shakespeare assistant helping a beginner understand the plays
Hamlet, Macbeth, and Romeo and Juliet. Use the retrieved passages below
to answer the question. You may reason across multiple passages.

When answering:
- Explain in plain English a beginner can understand
- Cite the play, act, and scene you drew from
- Do not invent details not supported by the passages
- Keep your answer concise and focused

If the passages are insufficient, say clearly what is missing and suggest
which play or act might contain the answer.

Retrieved passages:
{context}

Question: {question}
Answer:
\end{lstlisting}

\noindent\textbf{Baseline prompt} (\texttt{prompts/baseline\_prompt.txt}),
no retrieval:
\begin{lstlisting}
You are an assistant who will answer queries related to Shakespeare
whether or not the user knows about Shakespeare literature.

Answer the following question as best you can. If the question is
ambiguous, state your interpretation before answering. If unsure, say so
rather than inventing details.

Question: {question}
Answer:
\end{lstlisting}

\noindent\textbf{Stylised prompt} (\texttt{prompts/stylised\_prompt.txt}),
explicitly creative and capped at 150 words:
\begin{lstlisting}
Using the retrieved passages below as inspiration, respond in
Shakespearean Early Modern English style.

Rules:
- Use thee, thou, dost, hath, 'tis, wherefore where natural
- Maximum 150 words
- This is a creative response, not factual; do not present as evidence
- Write clearly enough that a beginner can still understand
- After the stylised response, add one plain-English sentence explaining it

Retrieved passages:
{context}

Question: {question}
Response (150 words max):
\end{lstlisting}

% ---------------------------------------------------------------
\section{Additional Output Examples}
\label{app:examples}

These worked examples show, for two instructor questions, the evidence
retrieved and the response produced by each system. They illustrate the
report's central finding: retrieval consistently locates the correct
scene, while the batch harness's stand-in generator
(Section~\ref{sec:evaluation}) fails to use it. Responses are quoted
verbatim and truncated; full records are in
\texttt{results/evaluation\_results.json}.

\noindent\textbf{Q04 --- Why does Macbeth kill Duncan?}
\begin{itemize}
  \item \emph{Baseline} (no retrieval): \textit{``Because Duncan was
  killed in the middle of a fight with Duncan\dots''} --- circular and
  factually empty; no evidence to draw on.
  \item \emph{Standard RAG} --- retrieved \emph{Macbeth}, Act~3 Scene~6;
  response: \textit{``If Duncan were the father of the prince, what would
  Duncan think?\dots''} --- relevant scene retrieved, but the generator
  did not ground its answer in it.
  \item \emph{Enhanced RAG} --- retrieved \emph{Macbeth}, Act~2 Scene~2
  (Lady Macbeth, the murder scene --- the correct evidence); response:
  \textit{``In the case of Duncan\dots he is a thief\dots''} --- the most
  relevant retrieval of the three, undone by generation.
\end{itemize}
\emph{Interpretation:} retrieval quality increases from baseline to
enhanced RAG (the enhanced system finds the actual murder scene), yet
answer quality does not, isolating generation as the bottleneck.

\noindent\textbf{Q02 --- What is the role of Lady Macbeth?}
\begin{itemize}
  \item \emph{Baseline}: \textit{``Lady Macbeth is a poet, poet, singer,
  dancer, musician\dots''} --- hallucinated, scored 1.
  \item \emph{Standard RAG} --- retrieved \emph{Macbeth}, Act~3 Scene~2;
  response: \textit{``Lady Macbeth is a servant, a servant\dots''} ---
  correct scene, degenerate (looping) generation.
  \item \emph{Enhanced RAG} --- retrieved a nearby Act~4 scene; response
  remained incoherent.
\end{itemize}
\emph{Interpretation:} a clear case of correct retrieval (the Macbeth
scenes) wasted by a generator too small to synthesise them --- the
failure mode the deployed Gemma~3~4B model is chosen to fix.

% ---------------------------------------------------------------
\section{GenAI Usage Log}
\label{app:genai-log}

\renewcommand{\arraystretch}{1.25}
\begin{longtable}{@{}p{3.4cm} p{6.2cm} p{6.2cm}@{}}
\caption{Consolidated GenAI usage log across all roles. Tool used:
Claude (Anthropic). Summaries are in our own words.}
\label{tab:genai-log}\\
\toprule
Prompt or task & Useful output & Verification / modification \\
\midrule
\endfirsthead
\multicolumn{3}{c}{\footnotesize\itshape Table~\ref{tab:genai-log} continued}\\
\toprule
Prompt or task & Useful output & Verification / modification \\
\midrule
\endhead
\bottomrule
\endlastfoot
Compare small models that run locally for a RAG pipeline &
Comparison of candidate models with parameter count, context window, and
RAM at 4-bit &
Cross-checked RAM against the Ollama library; used only to shortlist, not
copied into the report \\
Why Gemma~3~4B is suitable for this project &
Argument that a model with limited Shakespeare knowledge makes the
RAG-vs-baseline comparison genuine &
Confirmed by running the model with no context and seeing vague answers;
basis for the Section~\ref{sec:methodology} justification \\
Cross-encoder vs.\ bi-encoder: what does reranking improve? &
Bi-encoder encodes query and document separately; cross-encoder scores
the pair jointly for finer relevance &
Confirmed by observing the reranked pipeline surfaced different passages;
shaped \texttt{rag\_chatbot\_reranked.py} \\
How to choose retrieval depth $k$ &
Three-way trade-off of speed, relevance, and noise; suggested $k=5$ and
retrieve-15-rerank-to-5 &
Confirmed by running $k=3,5,8$ and seeing noise at $k=8$; kept $k=5$ \\
How embeddings handle archaic text &
Related modern and archaic terms occupy nearby regions; enrichment
bridges the vocabulary gap &
Confirmed by retrieving relevant passages lacking the query's exact words \\
Scene- vs.\ utterance-level chunking trade-offs &
Structured comparison of context preservation vs.\ retrieval noise and
corpus size &
Evaluated both against requirements; built both indices and compared
empirically \\
Effect of text enrichment on retrieval &
Prepending summaries and keywords improves query--passage alignment &
Verified by comparing retrieval with and without enrichment on the
instructor questions \\
Handling the 512-token embedding limit &
Recommended word-level overlap splitting to avoid silent truncation &
Implemented 400-word chunks with 50-word overlap; verified no truncation \\
Evaluation design and scoring rubric &
A 1--5 rubric over correctness, grounding, retrieval, usefulness, style &
Refined to match the specification; wrote our own 10 questions targeting
failure modes \\
LLM-as-judge for first-pass scoring &
How automated scoring works and where it is unreliable for small models &
Used as a first pass only; all scores reviewed manually against the CSV \\
Report structure and IEEE conventions &
Guidance on section ordering, technical depth, and citation style &
All content written and verified by us against the implementation \\
\end{longtable}

% ---------------------------------------------------------------
\section{Member Contributions}
\label{app:contributions}

This was a four-person project with the role distribution recommended in
the specification. All members contributed substantially; the summary
below gives equal weight to each role.

\renewcommand{\arraystretch}{1.25}
\begin{longtable}{@{}p{3.6cm} p{12.2cm}@{}}
\caption{Member roles and contributions.}
\label{tab:contributions}\\
\toprule
Role & Contribution \\
\midrule
\endfirsthead
\toprule
Role & Contribution \\
\midrule
\endhead
\bottomrule
\endlastfoot
\textbf{Data Engineering Lead} &
Acquired and inspected the provided corpus; performed cleaning
(stage-direction removal, whitespace normalisation, short-chunk
filtering); designed the processed record schema; implemented both
chunking strategies (229 scene-level chunks and 2{,}622 utterance-level
chunks) with modern-English summary/keyword enrichment; and built the
retrieval indices. Authored Section~\ref{sec:dataset}. \\
\textbf{Model Development Lead} &
Implemented the prompt-only baseline, the standard RAG pipeline, and the
enhanced reranked pipeline (cross-encoder + scene-level diversity);
designed the four prompt templates and the four interaction modes;
conducted the candidate-model comparison and wrote the model
justification; produced the system pipeline diagram. Authored
Sections~\ref{sec:methodology} and the implementation detail in
Section~\ref{sec:system}. \\
\textbf{Evaluation Lead} &
Designed the 15-question evaluation set (5 instructor, 10 group-designed)
and the 1--5 scoring rubric; implemented the automated evaluation harness
(\texttt{evaluate.py}) with LLM-as-judge first-pass scoring and manual
review; produced the comparative results, qualitative examples, and
failure analysis. Authored Sections~\ref{sec:evaluation}
and~\ref{sec:results}. \\
\textbf{Integration \& Reporting Lead} &
Designed and implemented the unified command-line interface
(\texttt{src/main.py}) with input validation, evidence display, stylised
labelling, graceful dependency handling, and a scripted demo mode;
organised the repository and documentation; assembled the four members'
sections into this single report and resolved cross-section consistency;
prepared the demonstration guide. Authored the Abstract,
Section~\ref{sec:intro}, the interface and repository content in
Section~\ref{sec:system}, Section~\ref{sec:genai},
Section~\ref{sec:conclusion}, and these appendices. \\
\end{longtable}

% ---------------------------------------------------------------
\section{Component Limitations and Integration Findings}
\label{app:limitations}

The following items were identified while integrating the four components
into one system and report. They are recorded here for transparency and
to direct future work; each is framed at the level of the artefact rather
than the author, and several were already flagged by the responsible
member.

\begin{enumerate}
  \item \textbf{Evaluation harness uses a stand-in generator
  (highest priority).} \texttt{evaluate.py} sets the generation model to
  \texttt{distilgpt2}, a placeholder, whereas the deployed and demonstrated
  system uses Gemma~3~4B. Consequently the correctness, grounding, and
  usefulness figures in Section~\ref{sec:results} and
  Appendix~\ref{app:eval-table} describe the placeholder, not the deployed
  model, and under-represent it. Retrieval-relevance scores are
  model-independent and are unaffected. \emph{Recommended action:} re-run
  the harness with the same Ollama call used by the interactive system and
  update the two result tables; the analysis structure is written to
  accommodate this without rewriting.

  \item \textbf{Two system framings.} The data and evaluation work describe
  systems as \emph{scene-level RAG} versus \emph{utterance-level enhanced
  RAG}, while the deployed system is framed as \emph{baseline / standard
  RAG / reranked RAG}. Both framings are valid, the project genuinely built
  scene and utterance indices, but a reader benefits from the explicit
  reconciliation now provided in Sections~\ref{sec:dataset}
  and~\ref{sec:evaluation}.

  \item \textbf{Generation, not retrieval, is the bottleneck.} Across the
  evaluation, retrieval relevance scored 5/5 while grounding scored near
  2/5. This is a genuine and useful finding, but in the current results it
  is entangled with item~1; separating a weak-placeholder effect from a
  true retrieval-to-generation gap requires the Gemma re-run.

  \item \textbf{Scene-level diversity cap (self-identified).} The reranked
  pipeline caps results at two per scene. As the Model Lead notes, on
  focused single-scene questions this can discard highly relevant passages;
  a minimum-coverage rule is the suggested replacement.

  \item \textbf{Minor repository cleanliness.} The repository contains a
  stray empty file (\texttt{results/koko}) that should be removed before
  submission. The dataset counts were re-verified during integration and
  reconcile correctly (3{,}613 raw utterances; 991 stage directions
  removed; 2{,}622 retained, matching the per-play totals).
\end{enumerate}

The original, unmodified report contributions of each member are retained
in the repository under \texttt{report/Varshini/}, \texttt{report/bobby/},
and \texttt{report/moinul/}; this integrated report assembles and
reconciles that work without altering the source files.
